\section{Experimental Setup}
\label{sec:setup}

\textbf{Protocol.} We follow the standard METR-LA protocol~\cite{li2018dcrnn,wu2019gwnet}:
12 input steps predict 12 output steps; windows are split chronologically 70/10/20
(23,974 / 3,425 / 6,850); and MAE, RMSE and MAPE are computed on raw mph over
\emph{non-zero} targets at horizons 3, 6 and 12 (15, 30 and 60\,min). The scaler,
weekly profile and sensor statistics are fit on the training period only. The test
split is used once, after all model selection.

\textbf{Inputs.} The Phase~1 deep models (LSTM, Graph WaveNet) receive two channels per sensor and step: z-scored speed
(mean 58.58, std 12.83, from observed training values) and time of day in $[0,1)$.
Following the reference pipeline, missing inputs keep their raw value 0, which
z-scores to $-4.57$; this is a distinctive value the network can learn to recognise.
ST-GAT and the hybrid instead use the masked, six-channel inputs of
Section~\ref{sec:dldata}. Tabular models receive the 34 features of
Section~\ref{sec:features}.

\textbf{Model selection and training.} Table~\ref{tab:hparams} lists the settings.
Consecutive windows share 11 of 12 steps, so the tabular models use every fourth
training window ($\approx$1.15\,M observed rows per horizon) with little loss of
information. The validation windows are split by parity: even windows select
hyper-parameters, odd windows drive early stopping. The Phase~1 deep models use
Adam~\cite{kingma2015adam} with gradient-norm clipping and keep the checkpoint with
the best validation MAE. Graph WaveNet uses the published architecture; we use
batch size 32 (instead of 64) to fit the 8\,GB unified memory of an Apple M3, and
train for at most 35 epochs (instead of 100) because of compute limits. All code
uses PyTorch~\cite{paszke2019} and scikit-learn~\cite{pedregosa2011}, with seed 0.

\begin{table}[t]
\centering
\caption{Hyper-parameters (selected on validation where a range is shown).}
\label{tab:hparams}
\setlength{\tabcolsep}{3pt}
\begin{tabular}{p{1.9cm}p{6.1cm}}
\toprule
Model & Settings \\
\midrule
Ridge & $\lambda\in\{10^{-2},1,10^{2},10^{3}\}$, standardised features, one model per horizon \\
Grad.\ boosting & absolute loss; $(\nu,\text{leaves})\in\{(0.1,31),(0.1,127),(0.05,63)\}$; min leaf 100; $\ell_2{=}1$; $\le$400 (800 for $\nu{=}0.05$) iterations, early stopping (patience 20); one model per horizon \\
LSTM & 2 layers, 64 units, dropout 0.1, MLP head, raw inputs, no graph; Adam lr $10^{-3}$, weight decay $10^{-4}$, batch 32, $\le$30 epochs of half the windows, patience 6 \\
Graph WaveNet & 4 blocks $\times$ 2 layers, kernel 2, channels 32/32/256/512 (residual/dilation/skip/end), diffusion order 2, embedding dim 10, dropout 0.3; Adam lr $10^{-3}$, weight decay $10^{-4}$, clip 5, batch 32, $\le$35 epochs, patience 8 \\
XGBoost per segment & one multi-output model per sensor, vector-leaf trees, absolute loss, depth 6, $\nu{=}0.1$, $\le$400 trees, early stopping (patience 30), every 2nd training window \\
ST-GAT & $d{=}64$, 2 ST blocks, $K_t{=}3$, 4+4 directional heads, node embedding 16, dropout 0.2, attention dropout 0.1; AdamW lr $10^{-3}$, weight decay $10^{-4}$, clip 5, batch 32, LR halving (patience 3), $\le$40 epochs of half the windows, patience 8 \\
Hybrid & ST-GAT as above + base-forecast embedding and fusion GAT, zero-initialised head; out-of-fold XGBoost base (5 blocked folds, purge 23) \\
Ablations & each ST-GAT / hybrid variant: exactly 12 epochs of half the windows, same seed \\
\bottomrule
\end{tabular}
\end{table}

\textbf{Fair comparison (Phase~2).} Every model, from Historical Average to the
hybrid, receives the same 12-step (1-hour) history and forecasts the same 12
horizons from the same anchors. The per-segment XGBoost uses 12 lags, as the ST-GAT
does. All models use the same chronological split, and all metrics share the same
missing-target mask. We report MAE, RMSE and MAPE at \emph{every} horizon,
separately for \emph{peak} periods (weekdays 06:00--10:00 and 15:00--19:00, by target
time) and off-peak, and for congested targets ($<$40\,mph). Pairwise differences are
tested with Diebold--Mariano tests, and overall MAE gets a 95\,\% day-block bootstrap
interval.

\textbf{Training.} Deep models are trained with the masked MAE of
Eq.~\eqref{eq:loss} in mph, batch size 32, AdamW (learning rate $10^{-3}$), and the
regularisers of Section~\ref{sec:dldata}. Main runs train for at most 40 epochs, each
on a random half of the training windows, with early stopping (patience 8) on
validation MAE. Ablations share one fixed budget of 12 such epochs and the same seed,
and the full model is re-trained under that budget so the rows are comparable. The
per-segment XGBoost models train on every second training window with early stopping
on every second validation window. All experiments ran on one Apple M3
laptop (8\,GB unified memory, Metal GPU) with PyTorch~\cite{paszke2019}, XGBoost
and scikit-learn~\cite{pedregosa2011}. The test split was used once, after all
design decisions.
